% Evans 11.5, printed pp.613--616; English PDF pp.630--633.
% Chinese start: printed p.433 / scan p.451.
\section{Problems}
\label{sec:11-5}
The following are abbreviated descriptions of the 17 source problems.
All solutions and additional qualifications are editorial reconstructions.

\subsection*{1. Generalized elastic system}
For $F=(-q(b),-p(a))$ with $p',q'>0$,
\[
 DF=\begin{pmatrix}0&-q'(b)\\-p'(a)&0\end{pmatrix},\qquad
 \lambda_\pm=\pm\sqrt{p'(a)q'(b)}.
\]
They are real and distinct. Corresponding vectors are
$(q'(b),-\lambda_\pm)$; substitution verifies the eigenvalue equation.

\subsection*{2. Shallow water and a change of conserved variables}
In $(h,q)$, $DF=\left(\begin{smallmatrix}0&1\\h-q^2/h^2&2q/h\end{smallmatrix}\right)$,
whose characteristic polynomial is $(\lambda-v)^2-h$.
For smooth $h>0$, expansion of $q=hv$ and use of
$h_t+(hv)_x=0$ give $v_t+vv_x+h_x=0$.
Hence the alternative smooth system has flux
$\widetilde F(h,v)=(hv,v^2/2+h)$ and matrix
$\left(\begin{smallmatrix}v&h\\1&v\end{smallmatrix}\right)$,
again with speeds $v\pm\sqrt h$.
This is equivalence for smooth solutions only. For example, the
original jump from $(h,v)=(1,0)$ to $(2,\sqrt3/2)$ has speed
$\sigma=\sqrt3$ and satisfies both original conservation laws.
For the alternative velocity law its flux jump is $11/8$, whereas
$\sigma[v]=3/2$; the jump fails. Nonlinear coordinate changes can
therefore alter admissible weak solutions even when both resulting
smooth systems happen to have conservation form.

\subsection*{3. Smooth matrices with rotating eigenspaces}
For $z\ne0$ put $\theta=1/z$ and
\[
 B(z)=e^{-1/z^2}\begin{pmatrix}\cos2\theta&\sin2\theta\\
 \sin2\theta&-\cos2\theta\end{pmatrix},\qquad B(0)=0.
\]
Every derivative is a finite sum of a bounded sine or cosine times
$z^{-N}e^{-1/z^2}$, which tends to zero. Induction proves $B\in C^\infty$.
The eigenvalues are $\pm e^{-1/z^2}$ for $z\ne0$, but coincide at zero.
The positive eigenline is spanned by $(\cos\theta,\sin\theta)$;
the negative line is its perpendicular. Along sequences $\theta=n\pi$
and $\theta=\pi/2+n\pi$, even the positive line has incompatible
limits. Sign choices cannot repair this, so no continuous unit
frame extends through zero. This failure concerns loss of strictness.

\subsection*{4. Shallow-water rarefaction curves}
In $(h,v)$ use $r_1=(1,-1/\sqrt h)$, $r_2=(1,1/\sqrt h)$.
Integrating $dv/dh=\mp h^{-1/2}$ gives
$2\sqrt h+v=C$ on $R_1$ and $2\sqrt h-v=C$ on $R_2$.
The gradient of each displayed invariant annihilates its own vector.

\subsection*{5. Shallow-water shock curves}
Write the base state as $(h_0,h_0v_0)$. Eliminating $\sigma$ from
$\sigma[h]=[hv]$, $\sigma[hv]=[hv^2+h^2/2]$ gives, for $h\ne h_0$,
\begin{equation}\label{eq:11-sw-shocks}
 v=v_0\pm(h-h_0)\sqrt{\frac{h+h_0}{2hh_0}},\qquad
 q=hv_0\pm(h-h_0)\sqrt{\frac{h(h+h_0)}{2h_0}}.
\end{equation}
One way to see the elimination is to introduce the relative mass flux
$j=h(v-\sigma)=h_0(v_0-\sigma)$; momentum gives
$j^2=(h^2-h_0^2)/(2(1/h_0-1/h))=hh_0(h+h_0)/2$.
The nontrivial branch speeds are
$\sigma=v_0\pm\sqrt{h(h+h_0)/(2h_0)}$.
If $h=h_0$, the mass jump forces $q=q_0$, so there is no missing
nonzero jump. For $q_0=0$, the formula is exactly
$q=\pm(h-h_0)h\sqrt{(1/h+1/h_0)/2}$.

\subsection*{6. Dam-break Riemann data with positive depths}
Let $h_l>h_r>0$ and both initial velocities be zero. The intermediate
depth $H$ is the unique root in $(h_r,h_l)$ of
\begin{equation}\label{eq:11-dambreak-root}
 2(\sqrt{h_l}-\sqrt H)
 =(H-h_r)\sqrt{\frac{H+h_r}{2Hh_r}}=:V.
\end{equation}
The left member strictly decreases and the right strictly increases:
for the latter square it and differentiate
$(H-h_r)^2(H+h_r)/(2Hh_r)$, or note both positive factors
$H-h_r$ and $\sqrt{1+h_r/H}$ have combined positive derivative.
The endpoint signs are opposite, proving existence and uniqueness.
Set $q_i=HV$, $\sigma=HV/(H-h_r)$. With $\xi=x/t$,
\[
 (h,v)=\begin{cases}
 (h_l,0),&\xi\le-\sqrt{h_l},\\
 ((2\sqrt{h_l}-\xi)^2/9,\ 2(\sqrt{h_l}+\xi)/3),
       &-\sqrt{h_l}<\xi<V-\sqrt H,\\
 (H,V),&V-\sqrt H\le\xi<\sigma,\\
 (h_r,0),&\xi>\sigma.
 \end{cases}
\]
The fan preserves $v+2\sqrt h=2\sqrt{h_l}$ and has
$\xi=v-\sqrt h$. Its right edge precedes the shock. Indeed
\[
 \sigma^2=\frac{H(H+h_r)}{2h_r}>h_r,\qquad
 (\sigma-V)^2=\frac{h_r(H+h_r)}{2H}<H.
\]
Thus $\sqrt{h_r}<\sigma<V+\sqrt H$, the Lax inequalities for a
2-shock; also $V-\sqrt H<\sigma$. The fan is continuous, the shock
satisfies both vector jumps, and the $O(t)$ disturbed interval proves
the initial trace. For all smooth entropy pairs the shock inequality
can also be checked without assuming small amplitude. Hold the left
state $(H,HV)$ fixed and follow its 2-shock branch toward smaller
$h$. Parameterize by $s=H-h>0$. Formula \eqref{eq:11-sw-shocks}
gives $\sigma(s)=V+\sqrt{h(h+H)/(2H)}$, strictly decreasing in $s$.
For $E(s)=[Q]-\sigma(s)[\eta]$, differentiated Rankine--Hugoniot gives
\[
 E'(s)=\sigma'(s)\{D\eta(u_r(s))\cdot(u_r(s)-u_l)
                      -\eta(u_r(s))+\eta(u_l)\}\le0
\]
by convexity, and $E(0)=0$. Thus this is an entropy solution in the
system sense, including for large positive depth ratios.
\PDEAnnotation{水坝问题不必小振幅}{局部 Riemann 定理只保证邻近状态。
本题通过单调求根、精确跳跃条件及沿激波支的熵耗散，
单独处理任意 $h_l>h_r>0$；不把局部定理扩大到真空或任意大数据。}

\subsection*{7. Barotropic invariants in the proper variables}
For $z=(\rho,q)$, use $r_1=(1,v-c)$ and $r_2=(1,v+c)$.
The gradients of \eqref{eq:11-gas-invariants} are
\[
 Dw_1=((c-v)/\rho,1/\rho),\qquad
 Dw_2=((c+v)/\rho,-1/\rho).
\]
Their respective products with $r_1,r_2$ vanish. This verifies the
source definition in conserved, rather than merely primitive, variables.

\subsection*{8. Compatibility for the alternative shallow-water system}
For $\widetilde F=(hv,v^2/2+h)$ the entropy relation reads
$Q_h=v\eta_h+\eta_v$, $Q_v=h\eta_h+v\eta_v$.
Equality of their cross derivatives gives
$\eta_{vv}=h\eta_{hh}$. This is necessary for a flux; convexity of
$\eta$ is a separate part of being an entropy pair.

\subsection*{9. Mechanical entropy for barotropic gas}
In conserved variables $(\rho,q)$ put
$\eta=q^2/(2\rho)+P(\rho)$ with $P''=p'/\rho$.
A compatible flux is
\[
 Q=\frac{q^3}{2\rho^2}+qP'(\rho).
\]
Direct differentiation gives $DQ=D\eta DF$. If the affine freedom
in $P$ is chosen so $\rho P'-P=p$, this is
$Q=v(\eta+p)$; without that normalization these expressions differ
by a constant multiple of the mass flux $q$ and one must use the
compatible expression above. In fact $\rho P'-P-p$ is a constant,
so choosing the additive constant in $P$ removes it.
The Hessian quadratic form is
\[
 D^2\eta[(a,b),(a,b)]
 =\frac{p'(\rho)}\rho a^2+\frac1\rho(b-va)^2>0
 \quad ((a,b)\ne0).
\]
Convexity here is in $(\rho,q)$, not asserted in $(\rho,v)$.

\subsection*{10. Gradient flux}
For $F=D\phi$ take $\eta=|z|^2/2$ and
$Q=z\cdot D\phi-\phi$. The derivative of $z\cdot D\phi$ is
$D\phi+z^TD^2\phi$; cancellation leaves $DQ=z^TD^2\phi=D\eta DF$.
The Hessian of $\eta$ is $I$.

\subsection*{11. Antigradient flux and the generalized elastic energy}
For $F=(\phi_2,\phi_1)$ and convex $\phi$,
$\eta=\phi$, $Q=\phi_1\phi_2$. Componentwise,
$D\eta DF=(\phi_1\phi_{12}+\phi_2\phi_{11},
\phi_1\phi_{22}+\phi_2\phi_{12})=DQ$.
For Problem 1 set $P'=p$, $R'=q$, and use
\[
 \eta(a,b)=P(a)+R(b),\qquad Q(a,b)=-p(a)q(b).
\]
Then $D^2\eta=\operatorname{diag}(p',q')>0$ and the compatibility
identity follows directly. This energy is obtained by calculation;
the flux itself need not come from a convex antigradient potential
with the same sign as in the first part.

\subsection*{12. Several space dimensions}
For $F=(F_1,\ldots,F_n)$ with $F_j:\Omega\to\R^m$, an entropy
pair is a convex $\eta$ and $Q=(Q_1,\ldots,Q_n)$ such that
$DQ_j=D\eta DF_j$ for each $j$.
Smooth solutions obey $\eta(u)_t+\operatorname{div}Q(u)=0$;
the entropy inequality for weak solutions is tested against every
nonnegative compactly supported scalar space-time function.

\subsection*{13. Nonlinear dielectric Maxwell energy}
With $B_t=-\operatorname{curl}E$, $D_t=\operatorname{curl}H$,
$E=D_D\eta$, $H=D_B\eta$, the chain rule gives
\[
 \eta_t=H\cdot(-\operatorname{curl}E)+E\cdot\operatorname{curl}H
       =-\operatorname{div}(E\times H).
\]
The identity $\operatorname{div}(E\times H)
=H\cdot\operatorname{curl}E-E\cdot\operatorname{curl}H$
fixes the sign. To verify the algebraic entropy compatibility, the
$j$th fluxes for the state $(B,D)$ have components
$F^{B_i}_j=\epsilon_{ijk}E_k$ and
$F^{D_i}_j=-\epsilon_{ijk}H_k$. Hence
$D\eta\,dF_j=H_i\epsilon_{ijk}dE_k-E_i\epsilon_{ijk}dH_k
=d(E\times H)_j$ (summation over $i,k$).
Convexity of the given energy supplies the entropy condition.

\subsection*{14. The entropy chord criterion for a scalar shock}
Let $a=u_l$, $b=u_r$, $a\ne b$, and
$\sigma=(f(b)-f(a))/(b-a)$. Write
$\ell(k)=f(a)+\sigma(k-a)$.
For $a<k<b$, the Kruzhkov jump defect is
\[
 [q_k]-\sigma[\eta_k]=2(\ell(k)-f(k)).
\]
For $b<k<a$, the defect is $2(f(k)-\ell(k))$.
Outside the interval it is zero by Rankine--Hugoniot.
Thus for an increasing jump entropy requires $f\ge\ell$ on $[a,b]$;
for a decreasing jump it requires $f\le\ell$ on $[b,a]$.
This also suffices for every smooth convex entropy: on a bounded
interval any such entropy is an affine function plus a positive
superposition of $|u-k|$, with weights $\eta''(k)/2$; integrate
the corresponding jump inequalities. The affine part contributes
zero by the vector jump identity. Diagram~\ref{fig:tikz-11.4}
shows the increasing-jump geometry. A uniformly convex $f$ lies
strictly below every nontrivial chord, so only decreasing shocks
are allowed.

\subsection*{15. Exact cubic-flux Riemann solution}
For $f(u)=u^3-u$, $a=-1$, seek a shock to a positive state $b$
whose speed equals the left edge of a rarefaction:
\[
 \frac{f(b)-f(-1)}{b+1}=b^2-b=3b^2-1=f'(b).
\]
The nontrivial root is $b=1/2$ and $\sigma=-1/4$.
The entropy solution is
\[
 u(x,t)=\begin{cases}
 -1,&x/t<-1/4,\\
 \sqrt{(x/t+1)/3},&-1/4<x/t<2,\\
 1,&x/t\ge2.
 \end{cases}
\]
On $[-1,1/2]$,
$f(u)-\ell(u)=(u+1)(u-1/2)^2\ge0$.
The shock satisfies the chord criterion and the fan is continuous
at $x/t=2$. Rankine--Hugoniot and the $O(t)$ trace argument prove
the weak initial problem. The local scalar contraction theorem
proves uniqueness in the bounded entropy class with strong local traces;
these Riemann states are not in $L^1(\R)$.

\subsection*{16. The perturbed sonic shock is curved}
Take $0<\lambda<1/4$. Label left characteristics by $\alpha<0$:
\[
 a(\alpha)=\begin{cases}-1+\lambda,&\alpha\le-1,\\
 -1-\lambda\alpha,&-1<\alpha<0,
 \end{cases}\qquad
 x=\alpha+(3a(\alpha)^2-1)t.
\]
Their Jacobian in the sloping part is $1-6\lambda a t>0$, so this
left region has no characteristic crossings. The initial shock from
Problem 15 now encounters increasing left states. Let its position
be $s(t)$, its left trace $a(t)<0$, and its right trace
$b(t)=-a(t)/2$. Rankine--Hugoniot and tangency give
\begin{equation}\label{eq:11-sonic-shock}
 s'=3a^2/4-1=f'(b),\qquad
 s=\alpha+(3a^2-1)t,\qquad a=-1-\lambda\alpha.
\end{equation}
Differentiating the second equation and substituting the first yields
\[
 \alpha'=-\frac{9a^2}{4(1-6\lambda a t)},\qquad
 a'=\frac{9\lambda a^2}{4(1-6\lambda a t)}>0.
\]
Thus $s''=(3/2)aa'<0$: the shock bends to the left relative to
its initial line. An explicit parameterization until it reaches the
left plateau is, with $r=-a\in[1-\lambda,1]$,
\[
 t(r)=\frac4{15\lambda}(r^{-8/3}-r^{-1}),\quad
 \alpha(r)=\frac{r-1}\lambda,\quad
 s(r)=\alpha(r)+(3r^2-1)t(r).
\]
Indeed the linear ODE for $t(a)$ has integrating factor $(-a)^{8/3}$.
The plateau is reached at $t_0=t(1-\lambda)$; after that
$a=-1+\lambda$ and the shock continues linearly with speed
$3(1-\lambda)^2/4-1$, matching its previous position and tangent.

The right-hand fan is no longer entirely centered at the origin.
For each emission time $0\le\theta\le\min(t,t_0)$, the tangent ray
\[
 x=s(\theta)+s'(\theta)(t-\theta),\qquad u=b(\theta)
\]
 fills the region from the shock to $x=-t/4$. Since
$\partial_\theta x=s''(\theta)(t-\theta)<0$ before $t_0$, these
rays do not cross in their domain and define the new fan uniquely.
For $-t/4<x<2t$ retain the original centered fan of Problem 15;
for $x>2t$ retain $u=1$. On the left use the noncrossing initial
characteristics up to the shock. The match at $x=-t/4$ has $u=1/2$.
Each shock jump $a\to-a/2$ obeys
$f(u)-\ell(u)=(u-a)(u+a/2)^2\ge0$ on that interval,
so it is admissible even though the fan derivative can be unbounded
on the sonic boundary. The solution is continuous away from the
shock, piecewise smooth away from fan edges, and its bounded fluxes
permit exhaustion by such regions in the weak identity. Equivalently
one may use the scalar characteristic construction and its continuous
fan limit. The initial transition region is bounded with an $O(t)$
perturbation near the origin, yielding the required local trace.
For a quantitative early-time check,
\[
 s(t)=-\tfrac14t-\tfrac{27}{16}\lambda t^2+O(t^3).
\]
Keeping the old fan $\sqrt{(x/t+1)/3}$ directly against a curved
shock would violate the entropy tangency condition. The emitted
rays are an essential part of the perturbation.

\subsection*{17. Smooth scalar contraction in several dimensions}
Let $w=u-\widehat u$ and define the smooth vector field
$b=\int_0^1F'(\widehat u+\theta w)\,d\theta$.
Then $w_t+\operatorname{div}(bw)=0$.
To avoid a discontinuous terminal adjoint, use
$\eta_\delta(w)=\sqrt{w^2+\delta^2}-\delta$.
The chain rule gives
\[
 \partial_t\eta_\delta(w)+\operatorname{div}(b\eta_\delta(w))
 =\{\eta_\delta(w)-w\eta_\delta'(w)\}\operatorname{div}b.
\]
The expression in braces has absolute value at most $\delta$ and
is zero where $w=0$. On a finite time strip with common compact
spatial support and smooth bounded derivatives, integrate and use
$\delta\int|\operatorname{div}b|$ on that compact set as an error.
It tends to zero. Dominated convergence gives
$\|w(t)\|_1=\|w(0)\|_1$, which is stronger than contraction as long
as both solutions remain smooth. Uniform compact support on a finite
strip follows from bounded speeds and initially compact support in
the smooth regime, or may be assumed explicitly in this calculation.
After shocks form the smooth renormalization argument no longer
applies, and only entropy contraction remains available.
